\documentclass[11pt]{article} 

\usepackage{pdfsync} 
\usepackage{url}
\usepackage{amsmath, amssymb}
\usepackage{hyperref}

\newcommand{\R}{\mathbb{R}}

\topmargin=0cm
\oddsidemargin0mm
\textheight23.5cm
\textwidth16cm
\headsep0mm
\headheight0mm
\parskip 2pt

\title{\bf 
	Sorbonne Université\\
	Implementation of finite element methods (MU5MAM30)\\[1cm]
	Some project proposals
}

\author{}
\date{}

\begin{document}
\maketitle

\section*{Instructions}

\begin{itemize}

\item Projects can be tackled in groups of at most three students (but different groups
can work independently on the same subject). 
\item Reports should be brief and to the
point, and ideally typeset in Latex or Markdown. They must be accompanied with source
code and build instructions, and either uploaded on the git repo in your branch
(say inside the "projects" directory) or sent to me by email, in the latter case 
in a single archive (zip or equivalent) with base name format 
\texttt{name\_surname-MU5MAM30-project}. 
\item When your project is finished or you have scheduled a deadline for it,  
take an appointment (30min slot, on an individual basis even if working in group) 
by sending me an e-mail listing your availabilities in the foreseen window for the 
presentation (select one to two weeks), I will then answer with a slot proposal.
\item
It is not needed to prepare slides for the presentation, I will 
have read your code and report in advance. You will expose me the difficulties 
and achievements during 5 to 10 minutes, and then I will ask you some specific questions 
related to your submission, or slightly around it, to evaluate your scientific autonomy.
\end{itemize}

Since most of the projects are extensions in different directions of
materials that we treated during the class, you {\it might} combine more
than one to build up your (multi)project.


\medskip
In case of guidance needed, for the choice or the realization of the project, 
contact me by e-mail or show-up in my office in 16-26 325. Proposing an
alternate project is also accepted, but you must send me the proposal for
agreement first.

\pagebreak

%%%%%%%%%%%%%%%%%%%%%%%%%%%%%%%%%%%%%%%%%%%%%%%%%%%%%%%%%%%%%%%%%%%%%%%%%%%%%%%%%%%%%%%%%%%%%%%%%%%%%%%
\section*{Project 1 : $\mathbb{P}_2$-Lagrange finite elements}

For this project you are asked to implement a $\mathbb{P}_2$-Lagrange FEM version
of what we have been doing during the tutorials with $\mathbb{P}_1$ elements 
only.\\
After computing the mass and stiffness terms for the reference element, which
shouldn't be too complicated, the main challenge is to devise well the 
``local to global'' numbering of the degrees of freedom, which are no longer 
simply the mesh vertices.\\
In terms of git repository, you could : 
\begin{itemize}
\item create files {\tt include/fem/P2.h} and {\tt /src/fem/P2.cpp} along the
existing $\mathbb{P}_1$ versions;
\item modify {\tt include/fem/mass.h} and {\tt include/fem/stiffness.h} to take
into account both FEM versions. Alternatively, replace them with 4 files 
{\tt mass$\_$Pi.h}  and {\tt stiffness$\_$Pi.h} with $i = 1$ and $i = 2.$;
\item modify {\tt include/FEM/poisson.h} and {\tt include/FEM/navier$\_$stokes.h}
(and of course also the {\tt src} versions) to add an argument describing  
the FEM type used for the solver. I'd suggest to do it in the constructor of 
the classes {\tt PoissonSolver} and {\tt NavierStokesSolver}, in addition to 
the already existing mesh argument;
\item finally add this FEM type choice in the test executables 
located in {\tt src/bin}, for example
as command line arguments of the {\tt main} functions.
\end{itemize}

\noindent{\it Suggestion if you wish to dig further :} In class we have seen that higher
order methods have better rates of convergence, as the mesh size tends to zero 
and at comparable total number of degrees of freedom, provided the original
solution itself 
(i.e. the one without FEM approximation) is sufficiently regular. You
might wish to test that claim numerically on a Poisson problem. I'd suggest to
work of a flat torus (i.e. a 2D planar grid with periodic boundary
conditions) rather than the sphere. First because in the case of the sphere 
there is also the effect of the mesh itself approximating the smooth manifold.
Second because you can then more easily devise right-hand-sides $f$ for which
you can compute the solution explicitly on paper. To exhibit the order of
convergence, use a plot of $L^2$ and $H^1$ errors with respect to mesh size, 
using logarithmic scales on the axis.

\pagebreak

%%%%%%%%%%%%%%%%%%%%%%%%%%%%%%%%%%%%%%%%%%%%%%%%%%%%%%%%%%%%%%%%%%%%%%%%%%%%%%%%%%%%%%%%%%%%%%%%%%%%%%%
\section*{Project 2 : Adding a boundary and boundary constraints}
The goal of this project is to adapt our solvers for cases of
2D surface meshes with a boundary, an corresponding appropriate 
boundary conditions.

In the case of the Poisson solver, you might implement non homogeneous
Dirichlet boundary conditions, since it is the most interesting case 
from the point of view of implementation. Identifying which vertices
of the mesh are on the boundary is a prerequisite. You may assume that
the mesh is oriented and manifold, so that a boundary edge is an edge 
with exactly one incident triangle, and a boundary vertex is simply a  
vertex belonging to a boundary edge.  

The case of the Navier-Stokes solver in vorticity/stream function framework 
requires some thinking at the theoretical level first, regarding which 
boundary conditions are physically relevant. Recall that 
the velocity $u$ of the fluid and the stream function $\Psi$ were related
by the equation  
$$
u = \nabla^\perp \Psi.
$$
A natural set-up in the presence of boundaries is to assume a 
non-penetrating condition, that is
$$
u \cdot n = 0 \text{ on } \partial \Omega,
$$
where here $n$ denotes the normal to the boundary. This can be rewritten as
$$
\nabla^\perp\Psi \cdot n = 0, \qquad \text{or equivalently} \qquad   \nabla \Psi
\cdot \tau = 0,
$$
where here $\tau$ is a tangent vector to the boundary. The latter implies that 
$\Psi$ is constant on each connected component of $\partial \Omega.$ In the project
you may assume that $\partial \Omega$ is connected (as is the case e.g. for
$\Omega$ a half-sphere surface mesh, hence only slightly modifying our class 
framework, or even more simply a 2D square grid). In that case, since $\Psi$ is 
only also globally defined up to a fixed constant (adding a constant to $\Psi$
does not modify the physical quantity $u$), you can assume that 
$\Psi \equiv 0$ on $\partial \Omega.$ Therefore, for the Poisson solver part of
the resolution, you will simply have
$$
\Delta \Psi = \omega \text{ on } \Omega,\qquad \Psi = 0 \text{ on } \partial
\Omega. 
$$

\noindent{\it Suggestion if you wish to dig further :} Try discussing the 
boundary conditions for $\omega$ that we are implicitly assuming here 
by not modifying the weak formulation for $\omega$ with respect to the
case without boundary. Other more physically relevant conditions might 
impose the value of the tangential part of the velocity (e.g. a 
no-slip condition corresponding to $u \cdot \tau = 0$, which in addition to
the impermeability condition yields $u = 0$ on $\partial \Omega$), or a 
so-called Navier type condition. Both of these are mathematically more 
challenging (non local) in the context of the vorticity-stream formulation.

\pagebreak

%%%%%%%%%%%%%%%%%%%%%%%%%%%%%%%%%%%%%%%%%%%%%%%%%%%%%%%%%%%%%%%%%%%%%%%%%%%%%%%%%%%%%%%%%%%%%%%%%%%%%%%
\section*{Project 3 : Adding a Coriolis force term}
The goal of this project is to adapt our Navier-Stokes solver on the sphere
in order to incorporate a Coriolis force term. 

\noindent 
For the report it is expected that you recall what the Coriolis 
force is and how it translates into the formulation of the 
Navier-Stokes equations, and specifically in our case in the context 
of the vorticity / stream function formulation.\\
The literature you might find online will most probably rely on using 
spherical coordinates, despite their singularity at the poles. In class though, 
we have noticed that the vorticity / stream function formulation can be
written intrinsically. The same is true for the Coriolis
term, and if you are facing difficulty to obtain it I'd suggest looking
at equation $(11)$ of 
\href{https://www.mathnet.ru/links/7186db239fce310100fce16f1e2395df/dan7557.pdf}{that}
nice 4 pages paper\footnote{It is written in russian but you only need to read
the math, which as you know is a universal language :-)} from the late eighties by Ilyin and
Filatov. Although it also uses spherical coordinates, you should recognize all
of our terms (in particular our transport term $T$ is called $J$ there), and the slight
modification due to the Coriolis force (called $\ell$). Pay attention that what is denoted by 
$\omega$ in the paper is not what we called vorticity, it is the magnitude of the
earth rotation speed. When testing, you may try to use numerical values corresponding 
to reality, or exaggerating it to strengthen its effect. Pay attention to the
scales and units though (our ``earth mesh'' has been rescaled to the unit $S^2$ 
sphere). Any simulation of yours (that includes devising an initial vorticity) 
that vaguely resembles a typical atmospheric circulation would be highly
appreciated ! 

\pagebreak
%%%%%%%%%%%%%%%%%%%%%%%%%%%%%%%%%%%%%%%%%%%%%%%%%%%%%%%%%%%%%%%%%%%%%%%%%%%%%%%%%%%%%%%%%%
\section*{Project 4 : Elimination tree and direct sparse Cholesky solver}

For this project you will write your own Cholesky factorization code  
$$
A = LL^T
$$
for SPD sparse matrices $A$ in the CSR format (or CSC if you prefer columns 
over rows).

You will then use it instead of the conjugate gradient for our Navier-Stokes
solver. The factorization will be done before the simulation starts, and
only forward and backward substitutions will be required at each time step. 
 
As we have seen in class one key step of this factorization in the framework 
of sparse matrices is to
determine in advance which 
matrix entries in $L$ might end-up being non zero (the fill-in), because these may be
much more than in $A$ and it is necessary or at least useful to ``prepare'' $L$ (reserve memory
etc) before filling it. This phase is usually called the {\it symbolic} phase, because 
the potential non zero entries are located and counted in each row, but are not
yet computed.

Since the theoretical background behind the process requires a bit of thinking, 
you are suggested to use as a reference the book {\it Algorithms for Sparse 
Linear Systems} by J. Scott \& M. Tuma (in open access, already uploaded in the
git repo inside the ``notes'' directory). Mostly relevant to us here are sections 
4.2, 4.3, (5.1) and 5.2. 
In particular you will learn the notion of {\it elimination tree}, and you 
will directly translate algorithms 4.2 and 4.3 (for the symbolic phase) and 
then $5.x$ (choose $x \in {5, 6, 7}$ most adapted to your set-up,  
for the actual factorization) into code. For the factorization step,
it is important to avoid computing elements of $L$ which were not marked as
potential non zeros in the symbolic phase, but this actual factorization phase 
will remain the most computationally expensive one.


\noindent{\it Suggestion :} 
Already after the symbolic phase, you may test using your code (e.g. computing the amount of
fill-in) the importance of the ordering of the unknowns, as we discussed in
class (see also section 8.4 of the book). For that purpose you can freely use 
the {\tt src/mesh/cube2.cpp} file that implements the cube mesh with nested 
dissection ordering ``by hand''.  For the application to Navier-Stokes, only
keep the nested dissection version for better performance.

\pagebreak
%%%%%%%%%%%%%%%%%%%%%%%%%%%%%%%%%%%%%%%%%%%%%%%%%%%%%%%%%%%%%%%%%%%%%%%%%%%%%%%%%%%%%%%%%%%%%%%%%%%%%%%%%%%%%
\section*{Project 5 : A (convex) nonlinear elliptic problem}

Let $\Omega$ be a bounded domain in $\R^2$ and let $f : \Omega \to \R.$ 
The area of the graph of $f$ is given by 
$$
A(f) := \int_\Omega \sqrt{1+ |\nabla f|^2}. 
$$
The graph of $f$ is said to be a minimal surface / minimal graph if
$$
A(f) \leq A(g) \text{ for any } g:\Omega \to \R \text{ s.t. } f=g \text{ on }
\partial \Omega,
$$
in other words if its area is minimal among all graphs assuming the same boundary curve.
A minimal graph satisfies the minimal surface equation 
$$
{\rm div}\left( \frac{\nabla f}{\sqrt{1 + |\nabla f|^2}}\right) = 0 \text{ on } \Omega,
$$
which is nothing but the Euler-Lagrange equation associated to $A(\cdot).$

Given a planar 2D triangular mesh $\Omega_h$, and piecewise affine 
function $f$ over that mesh (hence representable using a $\mathbb{P}_1$-Lagrange finite
element basis), present and code an algorithm that will compute/converge to the
(unique) minimal graph $f_*$ in that same finite element space and assuming
the same boundary values as $f$.\\
{\it Hint :} Note that the set of admissible candidates is an affine subspace of
the $\mathbb{P}_1$-Lagrange space over $\Omega_h$, and the functional $A(\cdot)$ is convex on that
set. Any reasonable gradient descent method should therefore converge.
 
\noindent{\it Suggestion if you wish to dig further :}
In nature minimal surfaces (e.g. soap films) need not be graphs. The mathematical 
set-up required to study 
the so-called Plateau's problem (finding a minimal surface given a boundary
curve) is therefore somewhat more complex, but the resulting Euler-Lagrange
equation is simpler since it amounts to the homogeneous linear Poisson equation
$\Delta f = 0.$ It comes with additional nonlinear constrains though, and
besides $f$ is no longer a scalar function (above the graph was $z = f(x,y)$, 
in the parametric formulation $f(x,y,z) = (f_1(x,y), f_2(x,y), f_3(x,y))$ is
a vector valued function. Solutions (they need no longer be unique) may still be searched 
for using a P1 finite element methods, you might for example try to reproduce the 
strategy explained in \href{https://eudml.org/doc/221286}{that} paper, in section 5.

\pagebreak
%%%%%%%%%%%%%%%%%%%%%%%%%%%%%%%%%%%%%%%%%%%%%%%%%%%%%%%%%%%%%%%%%%%%%%%%%%%%%%%%%%
\section*{Project 6 : Using existing software}

Understanding the math and the difficulties/challenges associated to the
implementation of finite element methods is important in order to later 
build on solid grounds. Since the subject is broad, it is also helpful if not 
important to be able to reuse (and potentially improve or make evolve) existing
code bases. For this last subject proposal, the goal is to reproduce the   
2D Navier-Stokes solver that we coded from scratch, by  
using third party libraries/software for each step (meshing, build-up of
matrices and solution of linear systems). You might for example choose one 
among the following well established open source projects :
\begin{itemize}
\item {\bf FreeFem} \url{https://freefem.org}\\
More a software than a library, it comes with its own scripting language 
leading to short single file scripts. In particular, you won't write any 
C/C++ code if choosing it.
\item {\bf FEniCS}  \url{https://fenicsproject.org}\\
Or it's DOLPHINx Python wrapper, if you also wish to avoid C/C++.
\item {\bf Deal.II} \url{https://www.dealii.org}\\
\end{itemize}
Demos and tutorials exist for all three, they might be easier to begin with
w.r.t. reading their documentation. The first one, FreeFem, is developed 
here at SU. 

\medskip

\noindent Another project direction, equally interesting, would be to only replace our linear
system solver(s) by state of the art ones. I'd strongly suggest getting access to
the {\bf MUMPS} solver \url{https://mumps-solver.org} (award wining and also
made in France)






\end{document}
